\section{Methodological Audit Framework}
\label{sec:methodology}

To establish whether complex metaphorical architectures are empirically necessary, we construct a rigorous, mathematically defined framework that isolates baseline initialization, regularization strength, and representation anisotropy. This section presents our classical parametric routing formulations, our maximum-entropy initialization strategy, the L2 regularized calibration protocol, and our covariance-injected anisotropy stress test.

\subsection{Classical Parametric Router Formulations}
Let $h_b^{(3)} \in \mathbb{R}^D$ represent the activation feature vector of sample $b$ extracted at Layer 3 (the frozen early boundary of the network), where $D = 192$ is the hidden dimension. The classical parametric router is parameterized as a single linear layer defined by weight matrix $W_g \in \mathbb{R}^{K \times D}$ and bias vector $b_g \in \mathbb{R}^K$, where $K = 4$ is the number of task-specific Low-Rank Adaptation (LoRA) experts (adapters) corresponding to MNIST, Fashion-MNIST, CIFAR-10, and SVHN, respectively.

We evaluate two gating mechanisms to map routing logits to ensembling weights $\boldsymbol{\alpha}_b \in \mathbb{R}^K$:
\begin{enumerate}
    \item \textbf{Softmax Router (Competitive):} Forces competitive, zero-sum space allocation across adapters. Let $z_{k, b} = \mathbf{w}_k^T h_b^{(3)} + b_{g, k}$ be the logit for expert $k$ on sample $b$ (where $\mathbf{w}_k^T$ is row $k$ of $W_g$). The blending coefficient $\alpha_{k, b} = \text{Softmax}_k( W_g h_b^{(3)} + b_g )$ is computed as:
    \begin{equation}
        \alpha_{k, b} = \frac{\exp(z_{k, b})}{\sum_{j=1}^K \exp(z_{j, b})}
    \end{equation}
    where $b_{g, k}$ is the $k$-th element of $b_g$.
    \item \textbf{Cooperative Independent Sigmoid Gating (Sigmoid Router):} Relaxes the zero-sum constraint, allowing multiple adapters to activate simultaneously if their features are complementary. Using the same logit $z_{k, b}$, the blending coefficient is computed as:
    \begin{equation}
        \alpha_{k, b} = \sigma\left( z_{k, b} \right) = \frac{1}{1 + \exp\left( - z_{k, b} \right)}
    \end{equation}
    where $\sigma(\cdot)$ is the standard logistic sigmoid function.
\end{enumerate}

\subsection{Maximum-Entropy Zero-Initialization (Standard Zero-Initialization)}
Standard neural network routing heads are commonly initialized using random uniform or normal distributions (e.g., Kaiming or Xavier initialization). Under extreme data scarcity (such as a tiny calibration support set), this random starting state introduces a severe initial routing bias, forcing the optimizer to struggle against random directionality. 

To resolve this confounding factor, we introduce a \emph{Maximum-Entropy Zero-Initialization} prior. To avoid terminological inflation, we note that this is mathematically identical to standard zero-initialization, where all routing weights and biases are initialized to exactly zero:
\begin{equation}
    W_g = \mathbf{0}, \quad b_g = \mathbf{0}
\end{equation}
We frame this simple setup through a maximum-entropy lens to highlight its physical and information-theoretic property: prior to optimization, the network resides in a state of maximum entropy, representing complete, unbiased uncertainty. Evaluating this zero-initialized state yields:
\begin{itemize}
    \item \textbf{Softmax Router:} $\alpha_{k, b} = \frac{1}{K}$ for all $k$, which is mathematically identical to static Uniform Merging, serving as an unbiased control.
    \item \textbf{Sigmoid Router:} $\alpha_{k, b} = 0.5$ for all $k$, providing an even, unwarped activation contribution from each task-specific adapter.
\end{itemize}

\subsection{Proper L2 Regularized Calibration (Standard Weight Decay)}
The parameters $\{W_g, b_g\}$ are optimized on a calibration dataset $\mathcal{D}_{\text{cal}} = \{(x_i, y_i)\}$ using Empirical Risk Minimization under an L2 Frobenius-norm complexity penalty. In standard optimization, this corresponds exactly to classical L2 regularization or weight decay. Let $f(x_i; W_g, b_g)$ represent the output logits of our 14-layer Sandbox network under the routed activation-space blend. The optimization objective is formulated as:
\begin{equation}
    \min_{W_g, b_g} \frac{1}{|\mathcal{D}_{\text{cal}}|} \sum_{(x_i, y_i) \in \mathcal{D}_{\text{cal}}} \mathcal{L}\left( f(x_i; W_g, b_g), y_i \right) + \lambda \|W_g\|_F^2
\end{equation}
where $\mathcal{L}$ is the multi-task classification cross-entropy loss, and $\lambda$ is the L2 weight decay hyperparameter. We employ the term "Proper L2 Regularized Calibration" because we sweep $\lambda \in \{0.0, 10^{-4}, 10^{-2}, 1.0\}$ across both small-sample and large-sample regimes to identify the mathematically correct regularization strength required to calibrate the routing parameters under varying data budgets.

\subsection{Anisotropy Stress Test via Covariance Injection}
Pre-trained foundation models exhibit severe representation anisotropy (the "representation cone"), where token features collapse into a narrow directional manifold. This geometric collapse increases feature correlation, rendering cosine-based nearest-centroid methods highly sensitive to representation overlap. To evaluate the robustness of our routers under these conditions, we introduce a parameterized \emph{Anisotropy Stress Test} by injecting a controlled covariance matrix $\Sigma$ into our synthetic task signatures $v_k$:
\begin{equation}
    v'_k = \Sigma^{1/2} v_k
\end{equation}
where the covariance matrix $\Sigma \in \mathbb{R}^{D \times D}$ is modeled as a Toeplitz matrix parameterized by an entanglement coefficient $\rho \in [0.0, 0.5]$:
\begin{equation}
    \Sigma_{i, j} = \rho^{|i - j|}
\end{equation}
As $\rho$ increases from $0.0$ (perfectly orthogonal task subspaces) to $0.5$ (severe representational entanglement), the features collapse into a highly anisotropic cone, simulating realistic foundation model conditions.

\subsection{Analytical Sandbox (ICS) Architecture}
We evaluate our methodologies inside a high-fidelity 14-layer PyTorch-based Analytical Coordinate Sandbox (ICS) that directly simulates representational flow and expert adaptation dynamics:
\begin{itemize}
    \item \textbf{Network Depth \& Dimension:} $L = 14$ layers; hidden dimension $D = 192$.
    \item \textbf{Frozen Boundary:} Layers $1$ to $3$ are early frozen layers. Feature extraction occurs at Layer 3: $h_b^{(3)} = v'_k + \epsilon_b$, where $\epsilon_b \sim \mathcal{N}(\mathbf{0}, \sigma_k^2 I)$ represents task-specific noise.
    \item \textbf{Dynamic Blending Layers:} For layers $l \in [4, 14]$, the activations are blended dynamically sample-wise via a coordinate-based attraction dynamical system:
    \begin{equation}
        h_b^{(l)} = h_b^{(l-1)} + \sum_{k=1}^K \alpha_{k, b} \gamma_V (v'_k - h_b^{(l-1)})
    \end{equation}
    where $\gamma_V = 0.05$ is the constant scaling factor (representing the LoRA expert scaling), and $v'_k$ represents the covariance-injected task signature. This iteratively pulls the blended representations toward the target task signatures.
    \item \textbf{Unified Distance-Based Classification Head:} The final representation $h_b^{(14)}$ is mapped to multi-task classification logits using a unified negative squared Euclidean distance classifier:
    \begin{equation}
        \text{logits}_{b, k} = - \| h_b^{(14)} - v'_k \|_2^2 + b_k
    \end{equation}
    where $b_k$ is the task-specific bias offset. This unified distance metric, along with the biases $b = [0.0, 0.0, -0.90, -2.30]$ and noise standard deviations $\sigma = [0.05, 0.15, 0.40, 1.20]$, is calibrated to perfectly reproduce standalone expert accuracies reported in modern benchmarks (MNIST: 100\%, FashionMNIST: 100\%, CIFAR-10: 92.40\%, SVHN: 22.80\%).
\end{itemize}

\subsection{Evaluation Metrics: Trajectory Jitter}
To evaluate the smoothness and stability of the ensembling weight trajectories across layers, we define the \emph{Trajectory Jitter} metric. For a given sample $b$ and a network of depth $L$ where dynamic blending occurs in layers $l \in [4, L]$, the jitter is the mean L2-norm of the difference in routing coefficients between adjacent layers:
\begin{equation}
    \text{Jitter} = \frac{1}{L - 4} \sum_{l=5}^L \|\boldsymbol{\alpha}_b^{(l)} - \boldsymbol{\alpha}_b^{(l-1)}\|_2
\end{equation}
where $\boldsymbol{\alpha}_b^{(l)} = [\alpha_{1, b}^{(l)}, \dots, \alpha_{K, b}^{(l)}]^T$ is the vector of blending coefficients computed at layer $l$. A low jitter value indicates a smooth, stable trajectory, whereas a high jitter indicates rapid oscillations in ensembling weights across layers.
